
\TNchapter[Agentic systems represent a shift from static automation to adaptive, goal-directed intelligence embedded across the enterprise.]{Continuous Evaluation Systems}
\begin{TNtwocol}
\section{Production Evaluation Pipeline Design }
\subsection{Introduction}
Production evaluation moves from deterministic testing to continuous, non-deterministic assessment in real-time. It requires handling scale and asynchronous processing without blocking user interactions.

\subsection{Core Architecture}
\begin{itemize}
    \item \textbf{Asynchronous Evaluation}: Decouples assessment from agent execution using event-driven architectures to ensure zero latency impact.
    \item \textbf{CI/CD Integration}: Automated hooks for regression detection and staged rollouts (Canary deployments).
    \item \textbf{Metric Collection}: Tracking Performance (latency, cost), Quality (accuracy, hallucinations), and User Feedback.
\end{itemize}

\end{TNtwocol}
\begin{TNtwocol}
\subsection{Deployment Strategies}
Detailed \textbf{Trace Capture} (using OpenTelemetry standards) enables granular analysis. \textbf{Monitoring} checks for quality degradation and drift against Service Level Objectives (SLOs). Using evaluation data to trigger automated retraining closes the improvement loop.
\end{TNtwocol}
\begin{TNtwocol}
\section{Sampling Strategies for Live Traffic }
\subsection{Introduction}
Evaluating every interaction is cost-prohibitive. Sampling strategies balance coverage with cost while maintaining statistical significance.

\subsection{Sampling Fundamentals}
The \textbf{Coverage-Cost Tradeoff} weighs the need for detecting issues against the expense of LLM-as-judge calls. \textbf{Risk-Based Sampling} prioritizes high-stakes or user-facing interactions over internal tool usage.

\subsection{Methodologies}
\begin{itemize}
    \item \textbf{Random Sampling}: unbiased baseline (Simple, Stratified, Reservoir).
    \item \textbf{Adaptive Sampling}: Increases rates when errors or low confidence are detected.
    \item \textbf{Stratified Sampling}: Ensures representation across Use Cases, User Segments, and Interaction Complexity to prevent blind spots.
\end{itemize}
\end{TNtwocol}
\begin{TNtwocol}
\section{Online vs. Offline Evaluation Integration}
\subsection{Introduction}
Integrating controlled Offline evaluation (reproducible but subject to distribution shift) with real-world Online evaluation (accurate distribution but lacks ground truth) is essential.

\subsection{Evaluation Paradigms}
\begin{enumerate}
    \item \textbf{Offline}: Best for regression testing and rapid iteration. Limited by determining "correctness" without real user context.
    \item \textbf{Online}: Best for drift detection and validating system integration. Challenges include attribution and lack of ground truth.
\end{enumerate}

\subsection{Integration Architecture}
\textbf{Bidirectional Feedback Loops} allow online failures to populate offline test sets, ensuring future tests reflect reality. \textbf{Staged Deployment} moves from development to staging (offline + synthetic load) and then to production monitoring.
\end{TNtwocol}
\begin{TNtwocol}
\section{Automated Quality Scoring Systems}
\subsection{Introduction}
Automated scoring uses LLM-as-judge frameworks and trajectory metrics to evaluate agents at scale, enabling rapid iteration and regression testing.

\subsection{LLM-as-Judge}
Frameworks like \textbf{GEval} allow criteria-based evaluation using natural language (e.g., "empathy," "compliance"). Effective rubrics need atomic criteria, explicit scoring bands, and reference examples.

\subsection{Evaluation Strategies}
\begin{itemize}
    \item \textbf{Reasoning Layer}: \texttt{PlanQualityMetric} and \texttt{PlanAdherenceMetric} assess strategy and execution faithfulness.
    \item \textbf{Action Layer}: \texttt{ToolCorrectness} and \texttt{ArgumentCorrectness} check if the right tool was called with valid parameters.
    \item \textbf{End-to-End}: \texttt{TaskCompletionMetric} and \texttt{Trajectory Evaluation} (Exact vs. Any-Order match) measure overall success.
\end{itemize}
\end{TNtwocol}
\begin{TNtwocol}
\section{Human-in-the-Loop Evaluation Integration }
\subsection{Introduction}
Humans provide ground truth for edge cases, high-stakes decisions, and nuanced quality dimensions where automation fails.

\subsection{Workflows}
\textbf{Annotation Queues} route interactions for review based on filters (e.g., negative feedback).
\begin{itemize}
    \item \textbf{Single-Run}: Detailed assessment against rubrics.
    \item \textbf{Pairwise Annotation}: Side-by-side comparison for A/B testing (e.g., model selection).
\end{itemize}

\subsection{Domain Expertise}
Effective evaluation requires SMEs (Clinicians, Lawyers) for domain-specific correctness, rather than generic annotators. Calibration rounds ensure inter-rater reliability.
\section{Feedback Loop Implementation }
\subsection{Introduction}
Feedback loops transform evaluation insights into improvements by channeling data back into development (datasets, prompts, models).

\subsection{Cyclical Development Model}
The lifecycle involves: Test Case Creation $\rightarrow$ Component Evaluation $\rightarrow$ Experimentation $\rightarrow$ Production Monitoring $\rightarrow$ Dataset Revision.

\subsection{Test Case Libraries}
Effective libraries focus on coverage over volume, evolving with real user variations and adversarial prompts harvested from production.

\subsection{User Feedback}
Integration includes \textbf{Explicit} signals (Thumbs Up/Down, Categorical Feedback) and \textbf{Implicit} signals (Retry rates, Abandonment) to validate automated metrics.
\section{Evaluation-Driven Retraining Triggers}
\subsection{Introduction}
Automated triggers determine when to retrain or fine-tune models based on performance drift, distinguishing signal from noise.

\subsection{Drift Detection}
\begin{itemize}
    \item \textbf{Input Distribution Drift}: Concept drift (new user needs) or Covariate shift (feature changes).
    \item \textbf{Performance Degradation}: Tracking gradual declines (staleness) or sudden drops (broken dependencies).
    \item \textbf{User Feedback Trends}: Declining ratings or rising complaint volumes.
\end{itemize}

\end{TNtwocol}
\subsection{Statistical Rigor}
Using \textbf{Rolling Baselines} and \textbf{Confidence Intervals} prevents over-reaction to natural variance. Contextual normalization accounts for time-based patterns.
